\subsection{Fixed-parameter FlexQP reconstruction}\label{sec:flex}
FlexQP addresses elastic quadratic optimization with a splitting formulation;
Deep FlexQP additionally uses a learned numerical-parameter policy~\cite{flex}.
The implemented reference reconstructs the fixed equations and Algorithm~1 from
the paper. It is neither the authors' released implementation nor a reproduction
of a learned policy. Trained weights, training data and a reproducible policy
implementation were unavailable within the comparison. The fixed reconstruction
is nevertheless executable on its stated domain and is evaluated against the
same independent lifted original objective.

The admitted model has positive finite-side linear prices and no finite hard
rows or quadratic soft costs. After converting finite sides to inequalities,
the elastic problem can be written as
\begin{equation}\label{eq:flex-model}
 \min_{x,s\geq0}\ \tfrac12x^\top P x+q^\top x+w^\top s
 \quad\text{subject to}\quad Gx-s\leq h,\qquad w>0.
\end{equation}
Lower interval sides become signed inequalities with their own prices. This
conversion preserves directionality, but it does not justify giving a hard
physical row a finite weight. The adapter consequently refuses hard rows rather
than choosing a numerically large weight and calling it exact hardness.

The fixed numerical parameters are $\sigma_x=10^{-6}$, $\sigma_s=1$, $\rho=1$
and relaxation $\alpha=1.6$. Eliminating the slack-related block produces the
Schur system
\begin{equation}\label{eq:flex-system}
 K=P+\sigma_x I+cG^\top G,\qquad
 c=(\sigma_s^{-1}+\rho^{-1})^{-1}.
\end{equation}
The positive proximal terms stabilize the numerical subproblem. They do not
modify the scientific objective in \eqref{eq:flex-model}; treating $\sigma_x$
as an objective regularizer would change the target optimizer. This distinction
is checked through original objective evaluation and stationarity, not through
a comparison of the stabilized subproblem's objective.

The implementation uses Python and SciPy iteration around its prepared factors.
It is less mature and has different language-level overhead from the compiled
C backends. Observed runtime differences therefore measure these concrete
implementations under the protocol, rather than an implementation-independent
lower bound on the FlexQP formulation. The reconstruction is not retuned after
examining the held-out results. In particular, the comparison does not choose a
different proximal parameter for each unfavorable test instance.

Algorithm~\ref{alg:flex} states the executable boundary without conflating
fixed and learned variants. The original dual of \eqref{eq:flex-model} has
finite-side bounds determined by the linear prices. The implementation checks
stationarity and the projected dual-box residual rather than accepting only a
small change between successive iterates. A stable iteration can converge to
the wrong fixed point if the sign or price mapping is incorrect; original-space
conditions are needed to distinguish that defect from a genuinely small step.

\begin{algorithm}[htbp]
\caption{Execute the fixed FlexQP equations on an admitted pure linear elastic problem, preserving original prices and independently testing the returned original-space optimum.}\label{alg:flex}
\begin{algorithmic}[1]
\Require Compatible pure linear elastic inequalities and fixed numerical parameters
\State Refuse finite hard rows, quadratic prices and nonpositive admitted side prices
\State Convert finite interval sides into $Gx-s\leq h$ with original linear prices
\State Prepare \eqref{eq:flex-system} using the fixed proximal parameters
\State Iterate the paper's fixed splitting and relaxation updates
\State Check original stationarity and projected dual-box residuals
\State Recover interval dual signs and compare the original objective to lifting
\end{algorithmic}
\end{algorithm}

The relationship between a finite elastic optimum and hard feasibility requires
care. If the original hard inequalities are feasible and suitable multiplier
bounds are known, a sufficiently large exact linear price can recover a hard
solution. That result does not make every finite price exact, nor does it make
the output safe for a plant whose immutable limits have been softened. The
application explicitly describes the objective being priced and separates it
from the original hard mask. This is a modeling and acceptance requirement,
not a criticism unique to one elastic algorithm.

Compatibility exclusions remain part of the scientific denominator. Symmetric
and asymmetric elastic families with nonzero quadratic side prices are outside
this pure linear reconstruction. They are marked incompatible before timing,
not counted as inaccurate solves or assigned zero durations. Conversely,
admitted asymmetric linear cases are genuinely executed and checked, so an
unavailable learned policy is not used to excuse omission of the fixed published
reference. The complete six-path original comparison can therefore expose
both semantic limits and implementation costs without pretending that all
methods solve every penalty family.
Within a prepared fixed-policy operation, the signed inequality matrix and
its numerical factor can be retained while admitted vectors change. A new
hard mask, finite-side pattern or price structure is not such a vector
change. The reconstruction's boundary checks distinguish these cases before
reuse. This is also why the complete-call timing is more informative than
an isolated triangular solve: it includes the work required to admit and
recover the actual scientific problem. The theoretical proximal regularizers
and the practical implementation overhead must both remain visible when
interpreting a comparison with a mature compiled backend.
